\subsection{The exact positive theta input for the cubic discriminant}
\label{sec:theta-input}

\paragraph{P-F46 (theta contraction, with its gluing and coefficient spaces).}
Use the cubic discriminant form $\phi$ and the period-ring map
$\mathcal P$ from P-F41. Let $L=II_{2,26}$, and let
$M\subset L$ be the primitive cubic-fourfold lattice with
signature $(2,20)$; this is the sign reversal of the
$(20,2)$ convention in Casalaina-Martin--Laza. Its orthogonal
complement is $K(-1)$, where $K=E_6$ is positive definite.
For an even lattice write $A_M=M^\vee/M$ and use the
discriminant quadratic form $(x,x)/2$ modulo $\mathbb Z$.
The unimodular gluing gives an isometry of finite quadratic
modules
\[
                 \iota:A_M\longrightarrow A_K.
\]
Both groups have order three. Write their classes as
$0,+,-$, with $q(\pm)=2/3$ modulo $\mathbb Z$, matching
the two nonzero classes by $\iota$. Define
\[
 \begin{split}
 \delta(\tau)&=q\prod_{j\ge1}(1-q^j)^{24},
                  \qquad q=e^{2\pi i\tau},\\
 \theta_{K+\lambda}(\tau)
       &=\sum_{x\in K+\lambda}q^{(x,x)/2},\\
 g(\tau)&=\frac1{\delta(\tau)}
       \sum_{\mu\in A_M}\theta_{K+\iota(\mu)}(\tau)e_\mu.
 \end{split}
 \tag{P-F46a}
\]
Here the Weil representation convention is
$\rho_M(T)e_\mu=e^{2\pi i q(\mu)}e_\mu$.
Then $g$ is a weakly holomorphic modular form of weight
$-9$ and type $\rho_M$, and its Borcherds lift satisfies
\[
       \Psi_M(g)|_{\mathbb B}\doteq\phi,
       \qquad
       \mathcal P\bigl(\Psi_M(g)|_{\mathbb B}\bigr)
                   \doteq\Delta_{\rm cub}.
 \tag{P-F46b}
\]
The notation $\doteq$ means equality up to a nonzero
constant; the scalar normalization of a Borcherds product
is not being fixed implicitly. Thus the output has
automorphic weight $48$, coefficient degree $80$, and
stretching degree $16$, exactly as in P-F41.

Every Fourier coefficient of every component of $g$ is
a nonnegative integer. More precisely, let $p_{24}(r)$
count partitions of $r$ with each part assigned one of
twenty-four colors, and put $p_{24}(r)=0$ for $r<0$.
Then
\[
 [q^t]g_\mu
   =\sum_{r\ge0}p_{24}(r)
       \#\{x\in K+\iota(\mu):(x,x)/2=t+1-r\}.
 \tag{P-F46c}
\]
All these sums are finite. The complete terms of exponent
at most zero are
\[
        (q^{-1}+96)e_0
                   +27q^{-1/3}(e_++e_-).
 \tag{P-F46d}
\]
In particular the constant coefficient giving the
Borcherds weight is $96$, so that $96/2=48$.

The input $g$ is not an ordinary square-integrable modular
input: its unregularized weight-$-9$ Petersson norm
diverges at the cusp. This assertion concerns the displayed
input itself, not all possible theta comparisons.

There is also a precise obstruction to a direct Hilbert
coefficient dictionary. Set $\nu_0=0$ and
$\nu_\pm=2/3$, and let
\[
                  c_\mu(N)=[q^{N+\nu_\mu}]g_\mu.
\]
For every fixed coset $\mu$, every integer $a\ge1$,
and every integer $b$, the sequence $c_\mu(as+b)$
cannot equal the actual cubic stretching multiplicity
$f_5(s)=\dim T^{(5)}_s$ for all sufficiently large $s$.
In fact it eventually exceeds $f_5(s)$. The same conclusion
holds for a positive rational affine reindexing on any
infinite arithmetic progression on which the exponents
belong to the component's Fourier lattice.

\paragraph{Proof of the exact contraction and the overlattice normalization.}
The primitive embedding with complement $E_6$ is the one
used in Casalaina-Martin--Laza, Proposition~7.5, printed
p.~25, in the construction already checked in P-F41.
Set $L'=M\oplus K(-1)$. Because $L$ is unimodular and
the two orthogonal factors are primitive, its gluing
subgroup
$I=L/L'\subset A_M\oplus A_{K(-1)}$
is the graph of the isometry $\iota:A_M\to A_K$.
One can see the sign explicitly: isotropy says
$q_M(\mu)-q_K(\iota\mu)=0$. There is therefore no
dual-Weil-representation sign left over in the final
$M$-valued theta series.

More explicitly, the projections of $I$ to the two
discriminant groups are injective by primitivity.
Unimodularity gives $|I|^2=|A_M||A_K|$, while
injectivity gives $|I|\le\min(|A_M|,|A_K|)$.
These two statements force $|I|=|A_M|=|A_K|$.
The determinant-three calculation for $K$ below
therefore makes all three orders equal to three,
and both projections are isomorphisms.
In the pullback convention of Ma, equation~(2.1),
pp.~5--6, the scalar vector $1\in\mathbb C[A_L]$
pulls back to
\[
                  \sum_{\mu\in A_M}
                      e_\mu\otimes\overline e_{\iota\mu}.
 \tag{P-F46e}
\]
Here $\overline e_\lambda$ is the basis vector for
$K(-1)$, identified with the dual of $e_\lambda$ for
$K$. Contracting this sum with $\Theta_K$ gives
exactly the sum in (P-F46a), with coefficient one
on each coset and no division by three or extra
overlattice sum. This is also Ma's Example~3.13,
printed p.~18, the unimodular specialization of
his Theorem~1.1 and equation~(1.1), p.~2.

That theorem applies: $M$ has rank $22$, and its
Witt index is at most two, strictly smaller than
$\operatorname{rank}M-2=20$. Its scalar input on
$L$ is $\delta^{-1}$, of weight $-12$, with
integral principal part and constant coefficient
$24$. The positive-definite rank-six theta series
has weight three and type $\rho_K\simeq\rho_M$;
Ma's theta-contraction formula consequently gives
the weight-$-9$ form $g$. This is the required
input weight $1-20/2$ for $M$.
The theorem identifies $\Psi_M(g)$, up to a
nonzero scalar, with the $E_6$ quasi-pullback of
$\Phi_{12}$. Restricting to the same ball as in
Casalaina-Martin--Laza gives $\phi$; applying
P-F41 gives (P-F46b). This proves an actual
source-to-output identity, not merely a similarity
between two theta constructions.

\paragraph{The coefficient counts and the complete principal part.}
The product identity
\[
             \delta^{-1}=q^{-1}
                    \sum_{r\ge0}p_{24}(r)q^r
\]
follows by expanding each factor $(1-q^j)^{-24}$:
it records an arbitrary number of parts of size $j$
in each of twenty-four colors. Multiplication by
the lattice sum proves (P-F46c). For a fixed exponent
$t$ only $0\le r\le t+1$ can contribute, and
positive definiteness makes each norm shell finite.
This proves both finiteness and nonnegative integrality.

Here is an elementary count of all the lattice terms
needed for (P-F46d), including the nonzero cosets.
Take
\[
 A_2=\{(a,b,c)\in\mathbb Z^3:a+b+c=0\},
 \qquad v=(2,-1,-1)/3,
 \qquad
 K=A_2^3+\mathbb Z(v,v,v).
 \tag{P-F46f}
\]
The gluing vector has norm two and integral pairings
with $A_2^3$, so this is an even integral lattice.
Its index over $A_2^3$ is three, giving determinant
$3^3/3^2=3$. To identify it with $E_6$ rather than
assume the identification, let $\alpha_i,\beta_i$
be simple roots in its three $A_2$ factors, chosen
so that $v_i=(2\alpha_i+\beta_i)/3$, and put
$u=v_1+v_2+v_3$.
The six roots
\[
             -u,\ \alpha_1,\beta_1,\alpha_2,\beta_2,\alpha_3
\]
have the $E_6$ Cartan matrix: the central root $-u$
has branches of lengths two, two, and one.
They generate $K$, since
$\beta_3=3u-2\alpha_1-\beta_1-2\alpha_2-\beta_2-2\alpha_3$.
This verifies the stated lattice and determinant convention.

The $A_2$ lattice has six vectors of norm two.
Each nonzero coset $\pm v+A_2$ has minimum norm
$2/3$, attained by exactly three vectors.
Indeed the vectors in $v+A_2$ have coordinates
$k_i-1/3$ with $k_i\in\mathbb Z$ and $\sum k_i=1$;
their squared norm is $\sum k_i^2-1/3$.
Its minimum is $2/3$, exactly when the $k_i$ are
a permutation of $(1,0,0)$. Negation proves the
other coset case.

The zero coset of $K$ is the union of $A_2^3$ and
the two classes $(\pm v+A_2)^3$. Its norm-two
vectors therefore number
\[
                         3\cdot6+2\cdot3^3=72.
\]
There are no additional possibilities: an $A_2$
factor of norm two already accounts for the full
norm in the first class, and each of the three
factors must have minimum norm in the other two.

For the dual quotient, identify
$(A_2^\vee/A_2)^3$ with $(\mathbb Z/3)^3$.
The gluing subgroup is the diagonal line
$H=\langle(1,1,1)\rangle$, and its orthogonal
subgroup consists of $(a,b,c)$ with $a+b+c=0$.
Thus $A_K=H^\perp/H$ has three elements.
Represent a nonzero class by $(1,-1,0)$.
Adding $H$ gives the three cyclic permutations
of this pattern. Each has minimum norm $4/3$,
with $3^2=9$ vectors at that norm; hence the
coset has exactly $27$ minimum vectors.
The negative class has the same count. This proves
\[
 \theta_K=1+72q+O(q^2),\qquad
 \theta_{K+\pm}=27q^{2/3}+O(q^{5/3}).
 \tag{P-F46g}
\]
Since $p_{24}(0)=1$ and $p_{24}(1)=24$,
(P-F46g) gives (P-F46d) and the weight $48$.
The norm and coset calculation also verifies
$q(\pm)=2/3$ in the convention used above.

\paragraph{A proved failure of direct coefficient reindexing.}
The zero vector in the zero coset, and one minimum vector
in either nonzero coset, show from (P-F46c) that for
every integer $N\ge0$,
\[
                         c_\mu(N)\ge p_{24}(N+1)
                                            \ge p(N+1),
 \tag{P-F46i}
\]
where $p$ is the ordinary partition function. The second
inequality assigns a fixed color to every part of an
ordinary partition; the first uses a lattice vector
of half-norm exactly $\nu_\mu$.

For an elementary lower bound on $p(r)$, put
$k=\lfloor\sqrt r\rfloor$ when $r\ge1$.
Every subset of $\{2,\ldots,k\}$ has sum at most
$k(k+1)/2-1\le r$. Turn it into a partition of $r$
by adding parts equal to one. Distinct subsets give
distinct partitions, since their parts greater than
one recover the subset. Hence
\[
                         p(r)\ge2^{\lfloor\sqrt r\rfloor-1}.
 \tag{P-F46j}
\]
On the actual plethysm side, $\dim\operatorname{Sym}^3
\mathbb C^5=35$, so the inclusion of invariants gives
\[
       f_5(s)\le
       \dim\operatorname{Sym}^{5s}(\mathbb C^{35})
                       =\binom{5s+34}{34}.
 \tag{P-F46k}
\]
For fixed $a>0$ and $b$, (P-F46i)--(P-F46j) grow at
least exponentially in $\sqrt s$, whereas (P-F46k)
is a polynomial of degree $34$. More explicitly,
the logarithm of the lower bound divided by $\sqrt s$
has positive lower limit, while the logarithm of the
upper bound divided by $\sqrt s$ tends to zero.
This proves the strict inequality for all sufficiently
large $s$. A compatible rational affine reindexing
becomes an integral affine one upon restricting $s$
to its fixed arithmetic progression; the right side
of (P-F46k) remains polynomial in the new parameter.
This proves the claimed obstruction without using
an asymptotic formula for partitions or assuming
an exact Hilbert dimension for the invariant ring.

\paragraph{The precise cusp obstruction to unchanged ordinary input.}
Uniformly for $x$ in a bounded real interval,
\[
                g_0(x+iy)=e^{-2\pi i x}e^{2\pi y}
                                      +96+O(e^{-2\pi y})
                         \quad(y\longrightarrow\infty).
\]
The expansion is uniform because the remaining power
series converges absolutely on every sufficiently
small closed $q$-disc. Thus for all large $y$,
$|g_0(x+iy)|\ge\tfrac12e^{2\pi y}$.
With the standard invariant norm on the Weil
representation, $\|g\|^2\ge|g_0|^2$.
Consequently the unregularized Petersson integral
over the cusp satisfies
\[
 \int_Y^\infty\int_{-1/2}^{1/2}
            \|g(x+iy)\|^2 y^{-9}\,\frac{dx\,dy}{y^2}
       \ \ge\ \frac14\int_Y^\infty e^{4\pi y}y^{-11}\,dy
       =\infty.
 \tag{P-F46h}
\]
In particular $g$ has exponential, not polynomial,
cusp growth. It cannot simply be substituted for an
ordinary square-integrable automorphic input while
retaining that input hypothesis. Borcherds's
regularization handles precisely a weakly holomorphic
input of this kind; the statement does not assert
that every theta lift of it diverges or that no
regularized comparison is possible.
\hfill$\square$

\paragraph{The same form is a unitary Borcherds lift.}
There is also an exact unitary interpretation, not merely
an analogy with an orthogonal construction. Equip $M$ with
a generator $\sigma$ of its cubic-cover deck group,
chosen to match the complex embedding of the period ball,
so that
$\sigma^2+\sigma+1=0$ and $\sigma$ preserves its integral
bilinear form $B$. Put $F=\mathbb Q(\zeta)$,
$\mathcal E=\mathbb Z[\zeta]$, where
$\zeta=e^{2\pi i/3}$ acts by $\sigma$. On this
$\mathcal E$-lattice define the form, linear in the first
variable,
\[
 h(x,y)=\frac{1-\zeta^2}{3}B(x,y)
           +\frac{\zeta-\zeta^2}{3}B(x,\sigma y).
 \tag{P-F46l}
\]
The identities $B(\sigma x,y)=-B(x,y)-B(x,\sigma y)$
and $B(\sigma x,\sigma y)=B(x,y)$ give
$h(\sigma x,y)=\zeta h(x,y)$ and
$h(y,x)=\overline{h(x,y)}$. Taking the field trace gives
$\operatorname{Tr}_{F/\mathbb Q}h(x,y)=B(x,y)$.
Also $B(x,\sigma x)=-B(x,x)/2$, so
$h(x,x)=B(x,x)/2\in\mathbb Z$.
Thus the Hermitian signature is $(1,10)$, and the
underlying quadratic lattice is exactly $M$, with no
rescaling of its quadratic form.

Write $\Lambda$ for this Hermitian lattice. For
$x,y\in\Lambda$ and $a\in\mathcal E$, the trace
of $\overline a h(x,y)$ is $B(x,ay)\in\mathbb Z$.
By trace duality this says
$h(x,y)\in\mathfrak D_F^{-1}$. It proves integrality
in Hofmann's convention, and the diagonal calculation
proves evenness. The same trace-duality argument proves
that its Hermitian dual $\Lambda'$ is exactly $M^\vee$:
if $v\in M^\vee$, apply integrality to $ax$ for every
$a\in\mathcal E$ and $x\in M$; the converse follows
by taking the trace. Consequently the finite quadratic
module is precisely $A_M$, with
$q(v)=h(v,v)=B(v,v)/2$, and Hofmann's Weil
representation in equation~(20), p.~19, is precisely
$\rho_M$, not its dual.

The isotropic hypotheses of Hofmann's Theorem~8.1
also hold. The displayed lattice in the proof of
Casalaina-Martin--Laza, Proposition~7.5, p.~25,
contains two hyperbolic planes; after sign reversal
this supplies a nonzero isotropic vector of $M$.
Since $\mathcal E$ is a principal ideal domain,
saturating its $F$-line in $\Lambda$ gives a primitive
isotropic vector $\ell\in\Lambda$.
Choose $y\in M\otimes\mathbb Q$ with
$a=h(\ell,y)\ne0$, and put
\[
              y'=y-\frac{h(y,y)}{2a}\ell.
\]
Direct expansion gives $h(y',y')=0$ and
$h(\ell,y')=a$. Clearing denominators gives an
isotropic $\ell'\in\Lambda\subseteq\Lambda'$ with
$h(\ell,\ell')\ne0$. These are exactly Assumption~2.1,
p.~6, and the hypotheses in Theorem~8.1, p.~24;
no unimodularity hypothesis from the special introductory
version of that theorem is being imposed.

Now $g$ has weight $1-10=-9$, the required Weil type,
and integral negative Fourier coefficients. Hofmann's
Theorem~8.1 therefore produces a unitary Borcherds
product $\Xi_g$, of weight $96/2=48$, initially for
the unitary discriminant kernel
$\Gamma^U_\Lambda\subset U(\Lambda)$, with a
finite-order multiplier. His definition in
equation~(28), p.~25, is
\[
        \Xi_g=\alpha^*\Psi_M(g)\doteq\phi,
        \qquad \mathcal P(\Xi_g)\doteq\Delta_{\rm cub}.
 \tag{P-F46m}
\]
Here $\alpha$ is the holomorphic embedding obtained
by viewing a positive Hermitian line as an oriented
positive quadratic two-plane, as in his Section~4,
pp.~13--15. This is the actual cubic-cover ball
embedding already used in P-F41 and (P-F46b).
Thus the source supplies the exact unitary lift name
for the same automorphic form. Its equality with
$\phi$, together with the group and multiplier check
in P-F41, extends this particular form to the exact
cubic period group with its stated trivial multiplier;
we do not identify that group with the discriminant
kernel. Theorem~8.1(4) further gives, after compatible
scalar normalization,
$\Xi_{r g}=\Xi_g^r$ for every positive integer $r$.
These are exactly the discriminant powers of output
degrees $(48r,80r,16r)$, rather than an assertion about
all graded pieces of the invariant ring.

\paragraph{What passes, and the remaining supplied-method comparison.}
The positive input spaces in (P-F46c) have a concrete
basis indexed by a lattice vector and a twenty-four-colored
partition. The established map sends their generating
series, through the regularized multiplicative lift,
to one algebra element, the cubic discriminant. The
coefficient $96$ determines its automorphic weight;
it is not identified with the dimension of a graded
plethysm multiplicity space. The powers of that one
element give the polynomial subalgebra from P-F41,
with degrees $(48r,80r,16r)$ for its $r$th power.

This sharpens the comparison with the supplied
\emph{Remark}, printed p.~4, item~8: the known
construction already contains a specific positive-definite
theta factor, but also the explicit inverse modular
discriminant, the coset contraction, and the regularized
multiplicative unitary Borcherds lift (P-F46m). Items~9--10 and~12 of that note
concern the norms, divisibility, and periods of its
iterated unitary lifts. No map identifying those
operations with (P-F46a)--(P-F46b) is provided by the
present argument. Such a comparison must account for
the displayed principal part and the divergent ordinary
input norm, as well as the multiplication and three
output gradings. This is a concrete remaining operation,
not an exclusion based on a group name.

In particular, the growth argument proves that the
positive Fourier coefficient spaces in (P-F46c) cannot
themselves be all the graded pieces of the invariant
ring under a positive affine reindexing. It does not
exclude subquotients, a different choice of spaces,
or a nonlinear transform. No basis
of those pieces, Hilbert-series identity with this
input, uniform positivity algorithm, or PH2/SH bound
is obtained. The theorem is an exact specialization of
a known quasi-pullback formula, with explicit coefficient
and normalization checks; no historical novelty is claimed.

The additional primary source is Shouhei Ma,
\emph{Quasi-pullback of Borcherds products},
Theorem~1.1 and equation~(1.1), p.~2; Weil conventions
and pullback, Section~2.1 and equation~(2.1), pp.~4--6;
theta contraction, equation~(3.1), p.~12; Example~3.13,
p.~18:
\url{https://arxiv.org/pdf/1801.08333}.

The unitary interpretation uses Eric Hofmann,
\emph{Borcherds Products on Unitary Groups},
Hermitian lattice conventions and Assumption~2.1, p.~6;
discriminant kernel, p.~7; underlying quadratic lattice
and holomorphic embedding, Section~4, pp.~13--15;
Weil representation, equation~(20), p.~19;
Theorem~8.1, p.~24; and its defining pullback,
equation~(28), p.~25:
\url{https://arxiv.org/pdf/1210.2542}.
